\scanpage{32}
donc en vertu de (1.1) on a
\[
\xi_k=0\quad\Longleftrightarrow\quad\xi_{k'}=0
\]
où $k'=i+j'=\sum_0^r(p-1)p^\alpha=p^{r+1}-1$,
donc on trouve que
\[
\xi_k=0\quad\Longleftrightarrow\quad\xi_{p^{r+1}-1}=0
\]
Pour la $=$ raison, on trouve, si $\widetilde k\deleted{\unsure{k_0}}\inserted{=i+\widetilde j}$ satisfait les $=$ conditions que $k$, savoir $c_{i,\unsure{k'|\widetilde k}-i}=0$
\[
\xi_{\widetilde k}=0\quad\Longleftrightarrow\quad\xi_{p^{r+1}-1}=0
\]
où on choisit $r$ tel que $p^{r+1}\geq\deleted{\unsure{k}}j,\widetilde j\ \unsure{(\geq i)}$.
En comparant, on trouve
\[
\xi_k=0\quad\Longleftrightarrow\quad\xi_{\widetilde k}=0
\]
i.e. si \uline{un} $\xi_k$ avec $k\geq N$, $c_{k,k-i}\not\equiv0\ (p)$ est nul, tous le sont, cqfd.

\uline{Fin de la démonstration}. Il \deleted{suffit de} \inserted{faut prouver} $\xi_k=0\quad\forall k\geq N$. On a
\[
c_{\ell,k}\xi_k=c_{\ell,k-i}\xi_{k+\ell}\qquad\forall\ell\geq0
\]
et il suffit de prouver qu'il existe $\ell\geq0$ tel que $c_{\ell,k}\not\equiv0\ (p)$, et de plus $\xi_{k+\ell}=0$, ce qui est le cas \deleted{si} en vertu de 1.1. si $c_{i,(k+\ell)-i}\not\equiv0\ (p)$. Posons encore $k-i=j$ i.e. $k=i+j$, on veut
\mspage{31}
